\medskip

Un étudiant mange tous les jours au restaurant universitaire. Ce restaurant propose des plats végétariens et des plats non végétariens.

\begin{itemize}[label=$\bullet~~$]
\item Lorsqu'un jour donné l'étudiant a choisi un plat végétarien, la probabilité qu'il choisisse un plat végétarien le lendemain est $0,9$.
\item Lorsqu'un jour donné l'étudiant a choisi un plat non végétarien, la probabilité qu'il choisisse un plat végétarien le lendemain est $0,7$.
\end{itemize}

\medskip

Pour tout entier naturel $n$, on note $V_n$, l'évènement \og l'étudiant a choisi un plat végétarien le $n$-ième jour \fg{} et $p_n$ la probabilité de $V_n$.

Le jour de la rentrée, l'étudiant a choisi le plat végétarien. On a donc $p_1 = 1$.



\begin{enumerate}
\item
	\begin{enumerate}
		\item 

Lorsqu'un jour donné l'étudiant a choisi un plat végétarien, la probabilité qu'il choisisse un plat végétarien le lendemain est $0,9$. Or $p_1=1$ donc $p_2=0,9$.

		\item 
$p_3= P\left (V_3\right )$. On représente la situation par un arbre pondéré.

\begin{center}
\medskip
\pstree[treemode=R,levelsep=3cm,nodesepA=4pt,nodesepB=4pt]{\TR{$V_1$}}
{\pstree[nodesepA=4pt]{\TR{$V_2~$}\taput{$0,9$}}
	{\TR{$V_{3}$}\taput{$0,9$}
	\TR{$\overline{V_{3}}$}\tbput{$0,1$}
	}
\pstree[nodesepA=4pt]{\TR{$\overline{V_2}$}\tbput{$0,1$}}
	{\TR{$V_{3}$}\taput{$0,7$}
	\TR{$\overline{V_{3}}$}\tbput{$0,3$}
	}
}
\end{center}

$\left \{ V_2\;,\; \overline{V_2} \right \}$ forme une partition donc, d'après la formule des probabilités totales:

$\begin{aligned}
p_3=P\left(V_3\right )
&  = P\left(V_2 \cap V_3\right ) +P\left(\overline{V_2} \cap V_3\right)
= P\left(V_2\right) \times P_{V_2}\left(V_3\right) + P\left (\overline{V_2} \right) \times P_{\overline{V_2} }\left(V_3\right)\\
& = 0,9\times 0,9+0,1\times 0,7 = 0,88
\end{aligned}$
	
\item Sachant que le $3\up{\text{e}}$ jour l'étudiant a choisi un plat végétarien,  la probabilité qu'il ait choisi un plat non végétarien le jour précédent est:

$P_{V_3}\left (\overline{V_2}\right )= \dfrac{P\left(\overline{V_2} \cap V_3\right)}{P\left ( V_3\right )}
= \dfrac{0,1\times 0,7}{0,88} = \dfrac{7}{88}\approx 0,08$ au centième près.


	\end{enumerate}
	
	
\item On complète l'arbre pondéré ci-dessous:

\begin{center}
\pstree[treemode=R,levelsep=3cm,nodesepA=0pt,nodesepB=4pt]{\TR{}}
{\pstree[nodesepA=4pt]{\TR{$V_n$}\taput{$p_n$}}
	{\TR{$V_{n+1}$}\taput{\blue $0,9$}
	\TR{$\overline{V_{n+1}}$} \tbput{\blue $0,1$}
	}
\pstree[nodesepA=4pt]{\TR{$\overline{V_n}$}\tbput{\blue $1-p_n$}}
	{\TR{$V_{n+1}$}\taput{\blue $0,7$}
	\TR{$\overline{V_{n+1}}$}\tbput{\blue $0,3$}
	}
}
\end{center}

\item
$\left \{ V_n\;,\; \overline{V_n} \right \}$ forme une partition donc, d'après la formule des probabilités totales:

$\begin{aligned}
p_{n+1}=P\left(V_{n+1}\right)
&  = P\left(V_n \cap V_{n+1}\right) +P\left (\overline{V_n} \cap V_{n+1}\right)
= P\left(V_n\right) \times P_{V_n}\left(V_{n+1}\right) + P\left(\overline{ V_n} \right) \times P_{\overline{V_n}}\left(V_{n+1}\right)\\
& = p_n\times 0,9+ \left (1-p_n\right )\times 0,7
= 0,9p_n + 0,7 - 0,7p_n = 0,2p_n + 0,7
\end{aligned}$

\item On souhaite disposer de la liste des premiers termes de la suite $(p_n)$ pour $n \geqslant 1$.

Pour cela, on utilise une fonction appelée \texttt{repas} programmée en langage Python dont on propose trois versions, indiquées ci-dessous.


\begin{center}
{\scriptsize
\begin{ttfamily}

\begin{tabular}{ccc}
\multicolumn{1}{c}{Programme 1}&\multicolumn{1}{c}{Programme 2}&\multicolumn{1}{c}{Programme 3}\\ 
\begin{tabular}{l| l|}\cline{2-2}
1&def repas(n):\\
2&~~p=1\\
3&~~L=[p]\\
4&~~for k in range(1,n):\\
5&~~~~p = 0.2*p+0.7\\
6&~~~~L.append(p)\\
7&~~return(L)\\ \cline{2-2}
\end{tabular}
&
\begin{tabular}{l| l|}\cline{2-2}
1&def repas(n):\\
2&~~p=1\\
3&~~L=[p]\\
4&~~for k in range(1,n+1):\\
5&~~~~p = 0.2*p+0.7\\
6&~~~~L.append(p)\\
7&~~return(L)\\ \cline{2-2}
\end{tabular}
&
\begin{tabular}{l| l|}\cline{2-2}
1&def repas(n):\\
2&~~p=1\\
3&~~L=[p]\\
4&~~for k in range(1,n):\\
5&~~~~p=0.2*p+0.7\\
6&~~~~L.append(p+1)\\
7&~~return(L)\\ \cline{2-2}
\end{tabular}\\
\end{tabular}
\end{ttfamily}
}
\end{center}

	\begin{enumerate}
		\item Le programme qui permet d'afficher les $n$ premiers termes de la suite $(p_n)$ est le \no 1.
		
En effet, le \no 2 donne $n+1$ termes, et le \no 3 donne des valeurs supérieures à 1.		
		
		\item Avec le programme \no 1, le résultat affiché pour $n = 5$ sera la liste $\left [p_1\;,\; p_2 \;,\; p_3 \;,\; p_4 \;,\; p_5 \strut \right  ]$.
		
$p_1=1$; $p_2=0,9$; $p_3 = 0,88$;
$p_4 = 0,2p_3 + 0,7 = 0,2\times 0,88 + 0,7 = 0,876$	 et\\
$p_5 = 0,2p_4 + 0,7 = 0,2\times 0,876 + 0,7 = \np{0,8752}$
		
Le résultat affiché par \texttt{repas(5)} sera donc $\left [1\;;\; 0,9 \;;\; 0,88 \;;\; 0,876 \;;\; \np{0,8752} \strut \right]$.
		
	\end{enumerate}
\item On va démontrer par récurrence que $p_n = 0,125 \times 0,2^{n - 1} + 0,875$, pour tout naturel $n \geqslant 1$.

\begin{list}{\textbullet}{}
\item \textbf{Initialisation}

Pour $n=1$, on a: $0,125 \times 0,2^{1 - 1} + 0,875 = 0,125 + 0,875 =1$.

Or $p_1=1$, donc la propriété est vraie pour $n=1$.

\item \textbf{Hérédité}

On suppose la propriété vraie au rang $n$ c'est-à-dire que $p_{n}=0,125 \times 0,2^{n - 1} + 0,875$; c'est l'hypothèse de récurrence.

$\begin{aligned}
p_{n+1}
& = 0,2p_n + 0,7
= 0,2 \left (0,125 \times 0,2^{n - 1} + 0,875\right )+0,7\\
& = 0,125 \times 0,2^{n } + 0,2\times  0,875+0,7
= 0,125 \times 0,2^{n } + 0,175+0,7\\
& = 0,125 \times 0,2^{n } + 0,875
\end{aligned}$

La propriété est donc vraie au rang $n+1$.

\item \textbf{Conclusion}

La propriété est vraie au rang 1 et elle est héréditaire pour tout $n\geqslant 1$; d'après le principe de récurrence, la propriété est vraie pour tout entier $n$ non nul.
\end{list}

On a donc  démontré par récurrence que $p_{n}= 0,125 \times 0,2^{n -1} + 0,875$ pour tout $n\geqslant 1$.

\item 
$-1<0,2<1$ donc $\ds\lim_{n\to +\infty} 0,2^{n-1}=0$ donc  $\ds\lim_{n\to +\infty} 0,125 \times0,2^{n-1} + 0,875 = 0,875$

La limite de la suite $(p_n)$ est donc $0,875$.
\end{enumerate}

\bigskip
